% Glossary terms

\newglossaryentry{NAT}{
    name={NAT},
    description={Network Address Translation. Una tecnica usata dai router per modificare gli indirizzi IP dei pacchetti in transito, permettendo a più dispositivi di condividere un singolo indirizzo IP pubblico}
}

\newglossaryentry{proxy}{
    name={Proxy},
    description={Un componente software o hardware che funge da intermediario tra un client e un server, inoltrando le richieste e le relative risposte}
}

\newglossaryentry{API}{
    name={API},
    description={Application Programming Interface. Un insieme di definizioni e protocolli utilizzati per costruire e integrare il software applicativo}
}

\newglossaryentry{UUID}{
    name={UUID},
    description={Universally Unique Identifier. Un identificativo standardizzato a 128 bit utilizzato nei sistemi informatici per identificare in modo univoco informazioni e risorse}
}

\newglossaryentry{JSON}{
    name={JSON},
    description={JavaScript Object Notation. Un formato di testo standard aperto, leggero e leggibile dall'uomo, utilizzato per lo scambio di dati tra client e server}
}

\newglossaryentry{WebAssembly}{
    name={WebAssembly},
    description={Un formato di istruzioni binario portabile, progettato per consentire l'esecuzione di codice ad alte prestazioni all'interno dei browser web}
}

\newglossaryentry{Appwrite}{
    name={Appwrite},
    description={Una piattaforma open-source per lo sviluppo di applicazioni web e mobile, che fornisce una serie di servizi backend come autenticazione, database, storage e funzioni serverless}
}

\newglossaryentry{outbound}{
    name={Outbound},
    description={Comunicazione di rete direzionata verso l'esterno di una rete locale LAN}
}

\newglossaryentry{inbound}{
    name={Inbound},
    description={Comunicazione di rete direzionata dall'esterno verso la parte interna di una rete LAN}
}

\newglossaryentry{LRU}{
    name={LRU},
    description={Least Recently Used, modalità di eviction di una memoria volatile, dove, una volta raggiunta la capacità massima, il dato aggiunto meno di recente viene espulso, in favore di uno nuovo}
}

\newglossaryentry{FIFO}{
    name={FIFO},
    description={First In First Out, modalità di gesione di una coda di dati, nella quale i primi elementi aggiunti sono anche i primi ad essere rimossi}
}

\newglossaryentry{Throughput}{
    name={Throughput},
    description={Capacità di trasmissione effettiva che viene usata in una certa unità di tempo}
}

\newglossaryentry{TTL}{
    name={TTL},
    description={Time to Live, intervallo di tempo della durata di validità di un dato prima della sua scadenza o eliminazione automatica. È solitamente espresso in secondi ed è utilizzato per gestire la persistenza temporanea di elementi come messaggi, sessioni o risorse di cache}
}
